\SGAProof{Proof}
By 5.6 and SGA 1, II 4.10, it is enough to prove that the
differentials of the functions
\[
\begin{gathered}
  a(d-1)-ec,\qquad
  b(d+1)-ec',\qquad
  b'(d+1)-e'c,\\
  e-2t(bc-ac'),\qquad
  e'-2t(b'c'-a'c),\qquad
  aa'-bb'-d
\end{gathered}
\]
at the points of the identity section of \(G\) are linearly
independent.  If the differential of a lower-case letter is denoted
by the corresponding capital letter, these differentials are
\[
  D,\qquad
  2B,\qquad
  2B',\qquad
  E+2tC',\qquad
  E'+2tC,\qquad
  A+A'-D,
\]
\footnote{Editor's note: \(A+A'+B\) has been corrected to
\(A+A'-D\).}
and they are indeed linearly independent modulo every
\((t-\lambda)\), \(\lambda\in k\).
\footnote{Editor's note: ``\((t-a)\), \(a\in A\)'' has been
corrected to ``\((t-\lambda)\), \(\lambda\in k\)''.}

\par\medskip\noindent
\textbf{\textit{\SGARefLink{sga3:xix:lemma:5:8}{Lemma 5.8}.}}---
\textit{For \(s\in S\), \(s\ne s_0\), the fiber \(G_s\) is connected
and semisimple.}
\SGARefTarget{sga3:xix:lemma:5:8}

\SGAProof{Proof}
Indeed, since \(s\ne s_0\), \(\mathfrak g(s)\) is isomorphic to the
Lie algebra of \(\mathrm{PGL}_{2,\kappa(s)}\), and, on the other hand,
\(G_s=(G_1)_s\).  The automorphism group of the Lie algebra of
\(\mathrm{PGL}_2\) over a field of characteristic \(0\) is known to
be \(\mathrm{PGL}_2\) itself, which is connected and semisimple.
\footnote{Editor's note: the preceding sentence has been slightly
modified.}

\par\medskip\noindent
\textbf{\textit{\SGARefLink{sga3:xix:lemma:5:9}{Lemma 5.9}.}}---
\textit{The fiber \(G_{s_0}\) is solvable and has two connected
components, of the following form:}
\[
\begin{aligned}
  G^0_{s_0}
  &=
  \left\{
  \begin{pmatrix}
    a&0&0\\
    0&a^{-1}&0\\
    c&c'&1
  \end{pmatrix}
  \right\},
  &
  G^-_{s_0}
  &=
  \left\{
  \begin{pmatrix}
    0&b&0\\
    b^{-1}&0&0\\
    c&c'&-1
  \end{pmatrix}
  \right\}.
\end{aligned}
\]
\SGARefTarget{sga3:xix:lemma:5:9}

\SGAProof{Proof}
We have \(e=e'=0\), since \(t=0\) at \(s_0\).  The equations
\textup{(1)}, \textup{(1')}, \(\ldots\), \textup{(5)}, and
\textup{(D)} can then be solved immediately.

\par\medskip\noindent
\SGAInlineRefTarget{sga3:xix:lemma:5-10}\textbf{\textit{\SGARefLink{sga3:xix:lemma:5:10}{Lemma 5.10}.}}---
\textit{The matrix}
\[
  w=
  \begin{pmatrix}
    0&1&0\\
    1&0&0\\
    0&0&-1
  \end{pmatrix}
\]
\textit{is a section of \(G\) over \(S\), such that
\(w(s_0)\in G^-_{s_0}\).}
\SGARefTarget{sga3:xix:lemma:5:10}

Let us now prove \SGARefLink{sga3:xix:proposition:5-5}{Proposition 5.5}.  Denote by \(G^0\) the union of the
identity components of the fibers of \(G\), that is, the complement
of \(G^-_{s_0}\).  Since \(G\) is smooth over \(S\) along the identity
section by \SGARefLink{sga3:xix:lemma:5-7}{Lemma 5.7}, \(G^0\) is an open subgroup of \(G\), smooth
over \(S\), by Exposé VI\textsubscript{B}, \SGARefLink{thm:sga3-VIB-3.10}{3.10}.  By translation,
\(G\) is also plainly smooth at the points of \(w(S)\); hence \(G\)
is smooth over \(S\).
\footnote{Editor's note: the preceding sentence has been modified,
and the reference to Exposé VI\textsubscript{B}, \SGARefLink{thm:sga3-VIB-3.10}{3.10} has been
added.}

\par\medskip\noindent
\textbf{\SGAInlineRefTarget{sga3:xix:paragraph:5-11}5.11.}---
Consider the morphism
\[
  \mathbb G_{m,S}\longrightarrow G^0,
  \qquad
  z\longmapsto
  \begin{pmatrix}
    z&0&0\\
    0&z^{-1}&0\\
    0&0&1
  \end{pmatrix}.
\]
It is a monomorphism defining a torus \(T\) of \(G^0\).  We claim
that
\[
  T=\operatorname{Centr}_G(T)
   =\operatorname{Centr}_{G^0}(T).
\]
\SGARefTarget{sga3:xix:paragraph:5:11}
